\documentclass[10pt,a4paper]{amsart}
\usepackage[margin=19mm]{geometry}
\usepackage{amsmath,amssymb,mathtools}
\usepackage[scr=rsfs,cal=euler]{mathalfa}
\usepackage{xfrac}
\usepackage[unicode,hidelinks,bookmarks=false]{hyperref}
\newcommand{\ie}{i.e.\ }
\newcommand{\cf}{cf.\ }
\newcommand{\resp}{respectively\ }
\newcommand{\Subsection}{\subsection}
\providecommand{\nobreakdash}{\nobreak\hbox{-}}
\setlength{\emergencystretch}{2em}
\multlinegap0pt
\usepackage[all]{xy}
\usepackage{tikz}
\usepackage{amssymb}
\usepackage[shortlabels]{enumitem}
\newtheorem{thm}{Theorem}
\newcommand{\CC}{\mathbb{C}}
\newcommand{\dbar}{\bar \partial}
\newcommand{\im}{\mathrm{Im} \,}
\renewcommand{\ker}{\mathrm{Ker} \,}
\title{Hodge theory for \\ twisted log-differential forms}
\author{Junyan Cao}
\date{Source: arXiv:2602.14826v1 (CC BY 4.0)}
\begin{document}
\maketitle

\noindent\emph{Source and licence.} Hodge theory for \\ twisted log-differential forms. Version arXiv:2602.14826v1 (CC BY 4.0), distributed by arXiv under the Creative Commons Attribution 4.0 International licence, http://creativecommons.org/licenses/by/4.0/, as recorded in the arXiv licence metadata for this exact version and shown on the abstract page https://arxiv.org/abs/2602.14826v1. Full licence notice: \texttt{licenses/arxiv-2602.14826-CC-BY-4.0.txt}.\par\smallskip

\noindent\textit{Editorial scope.} Counted unit: the article's thm extencla (source0.tex lines 2109--2118). The article's complete original proof of it is delivered with it (source0.tex lines 2120--2158, one \texttt{proof} environment of 39 lines). No mathematics is written, repaired or re-derived: the statement and the proof are the article's own lines, in the article's own notation, and no retained line is substituted or re-flowed.  No further source-local context is needed: every cross-reference of every retained line resolves inside the two delivered spans.  Nothing was omitted from any retained span, so the package carries no omission record.  No retained line contains a citation, so the wrapper carries no bibliography and no bibliography entry.  No retained line contains a figure environment, an image-inclusion command or drawing code, so no image is delivered and the package declares no asset dependency.  Editorial formatting only, in addition: an AMS \texttt{amsart} preamble that restates the article's own declarations and macros used by the retained lines, namely the environments \texttt{thm} on separate counters, together with the standard packages those lines need.  The retained environments print in this document as Theorem 1 in order of appearance, whereas the article prints the same items with the numbering of its own full document.  The complete native source of the article as distributed in the arXiv e-print for this exact version is bundled unchanged as \texttt{sources/194746/source0.tex}.
\begin{thm}\label{extencla}
	Let $X$ be a compact K\"ahler manifold and let $(L,\dbar_L)$ be the trivial line bundle.
	Let $\alpha$ be a smooth $\dbar$-closed $(0,1)$-form on $X$.
	Let $u$ be a smooth $\dbar$-closed $(p,q)$-form on $X$ which extends to first order in the direction $\alpha$.
	Then there exists a smooth family $(u_t)_{t\in (\CC,0)}$ of $(p,q)$-forms such that
	\begin{equation}\label{intr4}
		u_0 = u, \qquad u_t \in \ker \dbar_t.
	\end{equation}
	In other words, extension to first order implies extension to a genuine neighborhood.
\end{thm}

\begin{proof}
	We aim to solve the following system:
	\begin{equation}\label{simp46}
		\begin{cases}
			\dbar(u_1) + \alpha \wedge u = 0,\\
			\dbar(u_2) + \alpha \wedge u_1 = 0,\\
			\ldots\\
			\dbar(u_{i+1}) + \alpha \wedge u_i = 0,\\
			\ldots
		\end{cases}
	\end{equation}
	Note that $u$ extends to first order along $\alpha$ if and only if the first equation in \eqref{simp46} can be solved.
	Similarly, $u$ extends to arbitrary order if and only if all equations in \eqref{simp46} can be solved.
	
	Let $u_0$ be a harmonic representative of $u$. Since the first equation in \eqref{simp46} is solvable by assumption, we have
	\[
	\alpha \wedge u_0 \in \im \dbar \cap \ker \partial.
	\]
	By the $\partial\dbar$-lemma, there exists $u_1 \in \im \partial$ such that
	\[
	\dbar(u_1) + \alpha \wedge u_0 = 0.
	\]
	
	To solve the second equation, note that since $u_1 \in \im \partial$, we have $\alpha \wedge u_1 \in \im \partial$. Moreover,
	\[
	\dbar(\alpha \wedge u_1) = - \alpha \wedge \dbar u_1
	= \alpha \wedge \alpha \wedge u_0 = 0.
	\]
	Hence $\alpha \wedge u_1 \in \im \partial \cap \ker \dbar$.
	By the $\partial\dbar$-lemma again, there exists $u_2 \in \im \partial$ such that
	\[
	\dbar(u_2) + \alpha \wedge u_1 = 0.
	\]
	The same argument applies inductively to the remaining equations.
	
	Finally, if each $u_i$ is chosen as the minimal solution of \eqref{simp46}\footnote{It is then automatically $\partial$-exact by the Hodge decomposition theorem.},
	one can prove convergence of the series $\sum_{i=0}^\infty z^i u_i$ for $|z|\ll 1$.
	In particular, $u$ extends to a genuine neighborhood.
\end{proof}

\end{document}
